% Source: tex/chapters/ch05/06_テスト技法.tex
\subsubsection{経験ベース技法}
経験ベース技法は、仕様やコードだけではなく、テスト担当者の知識、直感、過去の障害、ドメイン理解を活用する。

\paragraph{エラー推測}
エラー推測は、過去の欠陥履歴や経験に基づき、「ここに欠陥がありそうだ」と考えてテストケースを設計する技法である。形式的な技法では見つけにくい問題を見つけることがある。

\paragraph{探索的テスト}
探索的テストは、学習、テスト設計、テスト実行を同時に進める活動である。事前にすべてのテストケースを決めるのではなく、観測した振る舞い、リスク、失敗の兆候に基づいて次のテストを考える。要求が詳細でないアジャイル開発などでよく使われる。

\paragraph{その他の経験ベース技法}
アドホックテストは、技能、直感、経験に基づいてテストする。モンキーテストはランダムな操作で停止や異常を狙う。ペアテストは二人でテストを行い、一人が操作し、一人が観察する。ゲーミフィケーションは、テスト作業をゲーム要素に変換して参加を促す。クイックテストは小さなテストスイートで重大問題を早く見つける。スモークテストは、ビルド後に基本機能が動くかを確認し、本格的なテストに進める状態かを判断する。
